%\title{University of Bristol Thesis Template}
\RequirePackage[l2tabu]{nag}		% Warns for incorrect (obsolete) LaTeX usage
%
%
% File: main.tex
% Author: Nowab Md. Aminul Haq
% Description: Contains the thesis template using memoir class,
% which is mainly based on book class but permits better control of 
% chapter styles for example. This template is an adaptation and 
% modification of Oscar's.
% 
% Memoir is a flexible class for typesetting poetry, fiction, 
% non-fiction and mathematical works as books, reports, articles or
% manuscripts. CTAN repository is found at:
% http://www.ctan.org/tex-archive/macros/latex/contrib/memoir/
%
%
% UoB guidelines for thesis presentation were found at:
% http://www.bris.ac.uk/esu/pg/pgrcop11-12topic.pdf#page=49
%
% UoB guidelines:
%
% The dissertation must be printed on A4 white paper. Paper up to A3 may be used
% for maps, plans, diagrams and illustrative material. Pages (apart from the
% preliminary pages) should normally be double-sided.
%
% Memoir class loads useful packages by default (see manual).
\documentclass[a4paper,11pt,leqno,openbib,oneside,french]{memoir} %add 'draft' to turn draft option on (see below)
%
%
% Adding metadata:
\usepackage{gensymb}
\usepackage{datetime}
\usepackage{ifpdf}
\usepackage{indentfirst}
\ifpdf
\pdfinfo{
	/Author (Author's name)
	/Title (PhD Thesis)
	/Keywords (One; Two;Three)
	/CreationDate (D:\pdfdate)
}
\fi
% When draft option is on. 
\ifdraftdoc 
\usepackage{draftwatermark}				%Sets watermarks up.
\SetWatermarkScale{0.3}
\SetWatermarkText{\bf Draft: \today}
\fi
%
% Declare figure/table as a subfloat.
\newsubfloat{figure}
\newsubfloat{table}
% Better page layout for A4 paper, see memoir manual.
\settrimmedsize{297mm}{210mm}{*}
\setlength{\trimtop}{0pt} 
\setlength{\trimedge}{\stockwidth} 
\addtolength{\trimedge}{-\paperwidth} 
\settypeblocksize{696.24pt}{451.44pt}{*} 
\setulmargins{1in}{*}{*} 
\setlrmargins{*}{*}{1} 
\setmarginnotes{17pt}{51pt}{\onelineskip} 
\setheadfoot{\onelineskip}{2\onelineskip} 
\setheaderspaces{*}{\onelineskip}{*} 
\checkandfixthelayout
%
\frenchspacing
% Font with math support: New Century Schoolbook
\usepackage{fouriernc}
\usepackage{}
\usepackage[T1]{fontenc}
%
% UoB guidelines:
%
% Text should be in double or 1.5 line spacing, and font size should be
% chosen to ensure clarity and legibility for the main text and for any
% quotations and footnotes. Margins should allow for eventual hard binding.
%
% Note: This is automatically set by memoir class. Nevertheless \OnehalfSpacing 
% enables double spacing but leaves single spaced for captions for instance. 
\OnehalfSpacing 
%
% Sets numbering division level
\setsecnumdepth{subsection} 
\maxsecnumdepth{subsubsection}
%
% Chapter style (taken and slightly modified from Lars Madsen Memoir Chapter 
% Styles document
\usepackage{calc,soul,fourier}
\makeatletter 
\newlength\dlf@normtxtw 
\setlength\dlf@normtxtw{\textwidth} 
\newsavebox{\feline@chapter} 
\newcommand\feline@chapter@marker[1][4cm]{%
	\sbox\feline@chapter{% 
		\resizebox{!}{#1}{\fboxsep=1pt%
			\colorbox{gray}{\color{white}\thechapter}% 
	}}%
	\rotatebox{90}{% 
		\resizebox{%
			\heightof{\usebox{\feline@chapter}}+\depthof{\usebox{\feline@chapter}}}% 
		{!}{\scshape\so\@chapapp}}\quad%
	\raisebox{\depthof{\usebox{\feline@chapter}}}{\usebox{\feline@chapter}}%
} 
\newcommand\feline@chm[1][4cm]{%
	\sbox\feline@chapter{\feline@chapter@marker[#1]}% 
	\makebox[0pt][c]{% aka \rlap
		\makebox[1cm][r]{\usebox\feline@chapter}%
}}
%\makechapterstyle{daleifmodif}{
\renewcommand\chapnamefont{\normalfont\Large\scshape\raggedleft\so} 
\renewcommand\chaptitlefont{\normalfont\Large\bfseries\scshape} 
%\renewcommand\chapternamenum{} \renewcommand\printchaptername{} 
%\renewcommand\printchapternum{\null\hfill\feline@chm[2.5cm]\par} 
\renewcommand\afterchapternum{\par\vskip\midchapskip} 
%\renewcommand\printchaptertitle[1]{\color{gray}\chaptitlefont\raggedleft ##1\par}
%} 
\makeatother 
\chapterstyle{daleifmodif}
%
% UoB guidelines:
%
% The pages should be numbered consecutively at the bottom centre of the
% page.
\makepagestyle{myvf} 
\makeoddfoot{myvf}{}{\thepage}{} 
\makeevenfoot{myvf}{}{\thepage}{} 
\makeheadrule{myvf}{\textwidth}{\normalrulethickness} 
\makeevenhead{myvf}{\small\textsc{\leftmark}}{}{} 
\makeoddhead{myvf}{}{}{\small\textsc{\rightmark}}
\pagestyle{myvf}
%
% Oscar's command (it works):
% Fills blank pages until next odd-numbered page. Used to emulate single-sided
% frontmatter. This will work for title, abstract and declaration. Though the
% contents sections will each start on an odd-numbered page they will
% spill over onto the even-numbered pages if extending beyond one page
% (hopefully, this is ok).
\newcommand{\clearemptydoublepage}{\newpage{\thispagestyle{empty}\cleardoublepage}}
%
%
% Creates indexes for Table of Contents, List of Figures, List of Tables and Index
\makeindex
% \printglossaries below creates a list of abbreviations. \gls and related
% commands are then used throughout the text, so that latex can automatically
% keep track of which abbreviations have already been defined in the text.
%
% The import command enables each chapter tex file to use relative paths when
% accessing supplementary files. For example, to include
% chapters/brewing/images/figure1.png from chapters/brewing/brewing.tex we can
% use
% \includegraphics{images/figure1}
% instead of
% \includegraphics{chapters/brewing/images/figure1}
\usepackage{import}



% Add other packages needed for chapters here. For example:
\usepackage{lipsum}					%Needed to create dummy text
\usepackage{amsfonts} 					%Calls Amer. Math. Soc. (AMS) fonts
\usepackage[centertags]{amsmath}			%Writes maths centred down
\usepackage{stmaryrd}					%New AMS symbols
\usepackage{amssymb}					%Calls AMS symbols
\usepackage{amsthm}	
\usepackage{amsmath}%Calls AMS theorem environment
\usepackage{newlfont}					%Helpful package for fonts and symbols
\usepackage{layouts}					%Layout diagrams
\usepackage{graphicx}					%Calls figure environment
\usepackage{longtable,rotating}			%Long tab environments including rotation. 
\usepackage[utf8]{inputenc}			%Needed to encode non-english characters 
%directly for mac
\usepackage{colortbl}					%Makes coloured tables
\usepackage{wasysym}					%More math symbols
\usepackage{mathrsfs}					%Even more math symbols
\usepackage{float}						%Helps to place figures, tables, etc. 
\usepackage{verbatim}					%Permits pre-formated text insertion
\usepackage{upgreek }					%Calls other kind of greek alphabet
\usepackage{latexsym}					%Extra symbols
\usepackage[square,numbers,
sort&compress]{natbib}		%Calls bibliography commands 
\usepackage{url}						%Supports url commands
\usepackage{etex}						%eTeX�s extended support for counters
\usepackage{fixltx2e}					%Eliminates some in felicities of the 
%original LaTeX kernel
\usepackage{babel}		%For languages characters and hyphenation
\usepackage{color}                    				%Creates coloured text and background
\usepackage[colorlinks=true,
allcolors=black]{hyperref}              %Creates hyperlinks in cross references
\usepackage{memhfixc}					%Must be used on memoir document 
%class after hyperref
\usepackage{enumerate}					%For enumeration counter
\usepackage{footnote}					%For footnotes
\usepackage{microtype}					%Makes pdf look better.
\usepackage{rotfloat}					%For rotating and float environments as tables, 
%figures, etc. 
\usepackage{alltt}						%LaTeX commands are not disabled in 
%verbatim-like environment
\usepackage[version=0.96]{pgf}			%PGF/TikZ is a tandem of languages for producing vector graphics from a 
\usepackage{tikz}	
%geometric/algebraic description.
\usetikzlibrary{arrows,shapes,snakes,
	automata,backgrounds,
	petri,topaths}				%To use diverse features from tikz		
%							
%Reduce widows  (the last line of a paragraph at the start of a page) and orphans 
% (the first line of paragraph at the end of a page)
\widowpenalty=1000
\clubpenalty=1000
%
% New command definitions for my thesis
%
\newcommand{\keywords}[1]{\par\noindent{\small{\bf Keywords:} #1}} %Defines keywords small section
\newcommand{\parcial}[2]{\frac{\partial#1}{\partial#2}}                             %Defines a partial operator
\newcommand{\vectorr}[1]{\mathbf{#1}}                                                        %Defines a bold vector
\newcommand{\vecol}[2]{\left(                                                                         %Defines a column vector
	\begin{array}{c} 
		\displaystyle#1 \\
		\displaystyle#2
	\end{array}\right)}
\newcommand{\mados}[4]{\left(                                                                       %Defines a 2x2 matrix
	\begin{array}{cc}
		\displaystyle#1 &\displaystyle #2 \\
		\displaystyle#3 & \displaystyle#4
	\end{array}\right)}
\newcommand{\pgftextcircled}[1]{                                                                    %Defines encircled text
	\setbox0=\hbox{#1}%
	\dimen0\wd0%
	\divide\dimen0 by 2%
	\begin{tikzpicture}[baseline=(a.base)]%
	\useasboundingbox (-\the\dimen0,0pt) rectangle (\the\dimen0,1pt);
	\node[circle,draw,outer sep=0pt,inner sep=0.1ex] (a) {#1};
	\end{tikzpicture}
}
\newcommand{\range}[1]{\textnormal{range }#1}                                             %Defines range operator
\newcommand{\innerp}[2]{\left\langle#1,#2\right\rangle}                                 %Defines inner product
\newcommand{\prom}[1]{\left\langle#1\right\rangle}                                         %Defines average operator
\newcommand{\tra}[1]{\textnormal{tra} \: #1}                                                       %Defines trace operator
\newcommand{\sign}[1]{\textnormal{sign\,}#1}                                                   %Defines sign operator
\newcommand{\sech}[1]{\textnormal{sech} #1}                                                  %Defines sech
\newcommand{\diag}[1]{\textnormal{diag} #1}                                                    %Defines diag operator
\newcommand{\arcsech}[1]{\textnormal{arcsech} #1}                                       %Defines arcsech
\newcommand{\arctanh}[1]{\textnormal{arctanh} #1}                                         %Defines arctanh
%Change tombstone symbol
\newcommand{\blackged}{\hfill$\blacksquare$}
\newcommand{\whiteged}{\hfill$\square$}
\newcounter{proofcount}
\renewenvironment{proof}[1][\proofname.]{\par
	\ifnum \theproofcount>0 \pushQED{\whiteged} \else \pushQED{\blackged} \fi%
	\refstepcounter{proofcount}
	\normalfont 
	\trivlist
	\item[\hskip\labelsep
	\itshape
	{\bf\em #1}]\ignorespaces
}{%
	\addtocounter{proofcount}{-1}
	\popQED\endtrivlist
}
%
%
% New definition of square root:
% it renames \sqrt as \oldsqrt
\let\oldsqrt\sqrt
% it defines the new \sqrt in terms of the old one
\def\sqrt{\mathpalette\DHLhksqrt}
\def\DHLhksqrt#1#2{%
	\setbox0=\hbox{$#1\oldsqrt{#2\,}$}\dimen0=\ht0
	\advance\dimen0-0.2\ht0
	\setbox2=\hbox{\vrule height\ht0 depth -\dimen0}%
	{\box0\lower0.4pt\box2}}
%
% My caption style
\newcommand{\mycaption}[2][\@empty]{
	\captionnamefont{\scshape} 
	\changecaptionwidth
	\captionwidth{0.9\linewidth}
	\captiondelim{.\:} 
	\indentcaption{0.75cm}
	\captionstyle[\centering]{}
	\setlength{\belowcaptionskip}{10pt}
	\ifx \@empty#1 \caption{#2}\else \caption[#1]{#2}
}
%
% My subcaption style
\newcommand{\mysubcaption}[2][\@empty]{
	\subcaptionsize{\small}
	\hangsubcaption
	\subcaptionlabelfont{\rmfamily}
	\sidecapstyle{\raggedright}
	\setlength{\belowcaptionskip}{10pt}
	\ifx \@empty#1 \subcaption{#2}\else \subcaption[#1]{#2}
}

%An initial of the very first character of the content
\usepackage{lettrine}
\newcommand{\initial}[1]{%
	\lettrine[lines=3,lhang=0.33,nindent=0em]{
		\color{gray}
		{\textsc{#1}}}{}}
%
% Theorem styles used in my thesis
%
\theoremstyle{plain}
\newtheorem{theo}{Theorem}[chapter]
\theoremstyle{plain}
\newtheorem{prop}{Proposition}[chapter]
\theoremstyle{plain}
\theoremstyle{definition}
\newtheorem{dfn}{Definition}[chapter]
\theoremstyle{plain}
\newtheorem{lema}{Lemma}[chapter]
\theoremstyle{plain}
\newtheorem{cor}{Corollary}[chapter]
\theoremstyle{plain}
\newtheorem{resu}{Result}[chapter]


%Added packages latyr nee
\newcommand{\itemcheck}{\item[\checkmark]}
%pour les formess rectangles et autres
\usepackage{tikz}
%pour les tables of contents in each part
\usepackage{minitoc}
%pour regrouper les lignnes d un tableau
\usepackage{multirow}
%pour inserer un pdf
\usepackage{pdfpages}

\begin{document}
	% UoB guidlines:
	%
	% Preliminary pages
	% 
	% The five preliminary pages must be the Title Page, Abstract, Dedication
	% and Acknowledgements, Author's Declaration and Table of Contents.
	% These should be single-sided.
	% 
	% Table of contents, list of tables and illustrative material
	% 
	% The table of contents must list, with page numbers, all chapters,
	% sections and subsections, the list of references, bibliography, list of
	% abbreviations and appendices. The list of tables and illustrations
	% should follow the table of contents, listing with page numbers the
	% tables, photographs, diagrams, etc., in the order in which they appear
	% in the text.
	% 
	\frontmatter
	\pagenumbering{roman}
	%
	%
% File: Title.tex
% Author: V?ctor Bre?a-Medina
% Description: Contains the title page
%
% UoB guidelines:
% 
% At the top of the title page, within the margins, the dissertation should give the title and, if 
% necessary, sub-title and volume number. If the dissertation is in a language other than English, the 
% title must be given in that language and in English. The full name of the author should be in the 
% centre of the page. At the bottom centre should be the words ?A dissertation submitted to the 
% University of Bristol in accordance with the requirements for award of the degree of ? in the 
% Faculty of ...?, with the name of the school and month and year of submission. The word count of 
% the dissertation (text only) should be entered at the bottom right-hand side of the page.
%
%
\begin{titlingpage}
	\begin{SingleSpace}
		\calccentering{\unitlength} 
		\begin{adjustwidth*}{\unitlength}{-\unitlength}
			%\vspace*{13mm}
			\begin{center}
				%\rule[0.5ex]{\linewidth}{2pt}\vspace*{-\baselineskip}\vspace*{3pt}
				%\rule[0.5ex]{\linewidth}{1pt}\\[\baselineskip]
				% {\Large \textit{Subtitle}}\\
				%\rule[0.5ex]{\linewidth}{1pt}\vspace*{-\baselineskip}\vspace{3pt}
				%\rule[0.5ex]{\linewidth}{2pt}\\
				{\large REPUBLIQUE DU SENEGAL }~\\[0.1cm]
				\includegraphics[width=0.05\textwidth]{Logos/senegal_flag.png}~\\[0.1cm]
				{\large UNIVERSITE CHEIKH ANTA DIOP }~\\[0.2cm]
				\hrule~\\[0.1cm]
				\includegraphics[width=0.15\textwidth]{Logos/logo_ucad.png}~\\[0.1cm]
				\hrule~\\[0.2cm]
				{\large Ecole Supérieure Polytechnique de Dakar }\\
				{\large Département Génie Informatique }\\
				\vspace{12mm}
				{\large MEMOIRE DE FIN DE CYCLE }~\\[0.1cm]
				{\large Pour l'obtention du : }~\\[0.1cm]
				{\large DIPLÔME D'INGENIEUR DE CONCEPTION EN INFORMATIQUE }~\\[0.1cm]
%				\hrule~\\[0.1cm]
				\vspace{12mm}
				\textbf{\huge Etude et mise en oeuvre d'une application conforme OWASP }\\[4mm]
				\vspace{3mm}
				\hspace{3mm}
				\begin{center}
					{\large \textbf{Lieu de stage :}}\\
					HubSo\\
				\end{center}
				\vspace{3mm}
				\hspace{3mm}
					\begin{flushleft}
						{\large \textbf{Présenté et soutenu par :}}\\
						Papa Latyr \textsc{Mbodj}\\
					\end{flushleft}
				\vspace{3mm}
				\hspace{3mm}
					\begin{center}
						{\large \textbf{\emph{Maîtres de stage :} }}\\
						{\large {Ahmed Tidiane CISSE}}\\
						{\large {Mouhamadou Mansour Sy SAMB}}\\
					\end{center}
				\vspace{3mm}
				\hspace{3mm}
					\begin{flushright}
						{\large \textbf{Professeur encadreur :} }\\
						{\large{M. Ibra DIOUM}}\\
					\end{flushright}
				\vspace{12mm}
				{\large Année universitaire : 2017-2018 }~\\[0.1cm]
				%\textsc
				%\clearpage
			\end{center}
			%\begin{flushright}
			%{\small Word count: ten thousand and four}
			%\end{flushright}
		\end{adjustwidth*}
	\end{SingleSpace}
\end{titlingpage}
\clearpage

	
	%
	%
% File: declaration.tex
% Author: V?ctor Bre?a-Medina
% Description: Contains the declaration page
%
% UoB guidelines:
%
% Author's declaration
%
% I declare that the work in this dissertation was carried out in accordance
% with the requirements of the University's Regulations and Code of Practice
% for Research Degree Programmes and that it has not been submitted for any
% other academic award. Except where indicated by specific reference in the
% text, the work is the candidate's own work. Work done in collaboration with,
% or with the assistance of, others, is indicated as such. Any views expressed
% in the dissertation are those of the author.
%
% SIGNED: .............................................................
% DATE:..........................
%
\chapter*{VISA}
\begin{SingleSpace}
	\begin{quote}
		\large {
			Je soussigné, Dr. Ibra Dioum, enseignant à l’École Supérieure
			Polytechnique de Dakar et encadrant, autorise M. Papa Latyr Mbodj, étudiant en cinquième année au Département Génie Informatique de l’École Supérieure
			Polytechnique de Dakar, option INFORMATIQUE à déposer son mémoire de fin de cycle DIC.} 
		\vspace{6cm}
		\begin{center}
			\large{Signature de l'encadrant}\\
			\vspace{12mm}
			\line(1,0){150}\\
			\large {Dr. Ibra Dioum}\\
			\vspace{2mm}
			\large {Département Génie Informatique}\\
			\large {École Supérieure Polytechnique}\\
			\large {Université Cheikh Anta Diop}\\
		\end{center}
		%
		
		
		%\end{center}
		%\noindent
		%\hspace{-0.75cm}\textsc{SIGNED: .................................................... DATE: ..........................................}
	\end{quote}
\end{SingleSpace}
%\clearpage
	
	%
	%
% file: dedication.tex
% author: V?ctor Bre?a-Medina
% description: Contains the text for thesis dedication
%

\chapter*{Dédicaces}

Je dédie ce mémoire à : \\
\begin{itemize}
	\itemcheck \textbf{mes chers parents Coumba MBODJ et Adama MBODJ} pour l’amour qu’ils m’ont porté ainsi que pour l’éducation qu’ils m’ont donnée. Je ne vous serai jamais assez reconnaissant. Que Dieu vous bénisse et vous accorde longue vie ;
	\itemcheck  \textbf{ma grand-mère Adja Fatou GUEYE, mon grand-père Ibrahima MBODJ et ma tante Fatou MBODJ}, vous êtes partis mais vous resterez à jamais dans mon cœur. Que le paradis soit votre demeure éternelle ;
	\itemcheck \textbf{ma tante Aissatou MBODJ} ;
	\itemcheck tous \textbf{mes oncles et tantes} ;
	\itemcheck tous \textbf{mes cousins et cousines} ;
	\itemcheck \textbf{toute la famille} ;
	\itemcheck \textbf{mon ami et frère de toujours, Aly SOW} ;
	\itemcheck \textbf{mon ami et frère de toujours, Thierno THIAM} ;
	\itemcheck \textbf{tous mes parrains et marraines particulièrement Papa Mamadou Ndiaye, Mansour Barro Diop et Mouhamed Samb} ;
	\itemcheck \textbf{tous les frères de «DoorWaar Team» et de la promo DUT 2014-2015} : Goss, Alfred, Thier, Cheikhna, Dave, Dio, Hamza, Maodo, Ouz Ndartoute, Bamba ainsi qu'à leurs familles ;
	\itemcheck \textbf{toutes mes sœurs de la promo DUT 2014-2015} ;
	\itemcheck \textbf{tous les amis} ;
	\itemcheck \textbf{toutes les personnes} qui, de près ou de loin, ont contribué à l’aboutissement de ma formation et à faire de moi celui que je suis.
\end{itemize}
\clearpage
	
	%
	%
% file: dedication.tex
% author: V?ctor Bre?a-Medina
% description: Contains the text for thesis dedication
%

\chapter*{Remerciements}
%\begin{SingleSpace}

%\end{SingleSpace}
\clearpage
	
	%
	%
% File: declaration.tex
% Author: V?ctor Bre?a-Medina
% Description: Contains the declaration page
%
% UoB guidelines:
%
% Author's declaration
%
% I declare that the work in this dissertation was carried out in accordance
% with the requirements of the University's Regulations and Code of Practice
% for Research Degree Programmes and that it has not been submitted for any
% other academic award. Except where indicated by specific reference in the
% text, the work is the candidate's own work. Work done in collaboration with,
% or with the assistance of, others, is indicated as such. Any views expressed
% in the dissertation are those of the author.
%
% SIGNED: .............................................................
% DATE:..........................
%
\chapter*{Avant-propos}
\begin{SingleSpace}



\end{SingleSpace}
%\clearpage
	
	%
	
\chapter*{Résumé}
Balancing of an inverted pendulum is a classical problem in the field of Control Theory and Engineering. It's balancing is always challenging for the beginners in control engineering. This thesis deals with stabilization and control of both the linear and rotary inverted pendulum  systems. Simultaneous approach is applied to system analysis as well as controller synthesis. Great similarity of both the systems is pointed out during the derivation of the equations of motion using the Euler-Lagrange equation. Linear inverted pendulum system is analyzed in two experimental setups , pendulum on cart and two wheel self balancing vehicle. PID controller is used to control the system by trial and error based tuning. Potentiometer (pot) and gyroscopic sensor is used as feedback sensor of the controlled system. Rotary inverted pendulum system is analyzed in a experimental setup developed by $QUANSER$ and is balanced by Pole placement method. Rotary encoder is used as the feedback sensor of the controlled system. Later, Linear Quadratic Regulator (LQR) is designed for optimum control of the pendulum. System response and controller gains are simulated in $MATLAB$ environment and after applying controller in experimental prototype, actual response of the system is reported. \\ 
\\ \textbf{Keywords}: Inverted pendulum,  PID controller, Pole placement method, LQR.
\clearpage

	
	
	%
	
\chapter*{Abstract}

Implementing secure Web and Mobile applications is a classic problem for many businesses. The rapid growth of connectivity, combined with the dramatic increase in the value of the data handled by these applications as well as the increasing use of new protocols and technologies have resulted in applications exposing both the organizations that are putting them in place as well as their users at significant risks.\\
Although the term « secure » is relatively relative, the fact remains that a certain degree of security is essential for these applications. To achieve this, an ideal approach is to consider security as an integral part of the project throughout its life cycle. Again, the parties should be aware of this and the resources to achieve this goal should be available.\\
This document is a late cycle memoir for obtaining the Design Engineering degree in Computer
Science. This presents the realization of a project carried out as part of an internship at HubSo. This project consists of the study and implementation of OWASP compliants Java EE web application and Android Mobile application.\\
We have designed a security features library that can be integrated into web and mobile applications. Our solution is based on a pre-existing open source API developed by OWASP, OWASP ESAPI. In addition, the library has been integrated with an existing application.\\
The library has been developed in Java SE and a developer guide has been proposed along with an OWASP compliant application architecture\\ \\
\textbf{Keywords}: Security, OWASP, Java

\clearpage

	
	%
	%\chapter*{Sigles et Abbréviations}
\addcontentsline{toc}{chapter}{Sigles et Abbréviations}
\markboth{chapter name}{Sigles et Abbréviations}
\printnomenclature[1in]
\clearpage
	
	
	%
	\renewcommand{\contentsname}{Table des matières}
	\maxtocdepth{subsubsection}
	\tableofcontents*
	\addtocontents{toc}{\par\nobreak \mbox{}\hfill{\bf Page}\par\nobreak}
	\clearpage
	%
	\listoftables
	\addtocontents{lot}{\par\nobreak\textbf{{\scshape Table} \hfill Page}\par\nobreak}
	\clearpage
	%
	\listoffigures
	\addtocontents{lof}{\par\nobreak\textbf{{\scshape Figure} \hfill Page}\par\nobreak}
	\clearpage
	
	%Introduction
	%\addcontentsline{toc}{chapter}{Introduction}
\chapter*{Introduction}
\lipsum[1-9]
\clearpage
	%
	
	%
	%
	% The bulk of the document is delegated to these chapter files in
	% subdirectories.
	\mainmatter
	%
	\import{misc/introduction/}{introduction}
	\import{parts/part01/}{part01.tex} % Présentation Générale
	\import{parts/part02/}{part02.tex} % Etat de l'art sur la sécurite informatique
%	\import{parts/part03/}{part03.tex} % Etat de l'art sur la securite des applications Web
	\import{parts/part04/}{part04.tex} % Analyse et Conception
	\import{parts/part05/}{part05.tex} % Réalisation
	\import{parts/part06/}{part06.tex} % Bilan et Perspectives
	\import{misc/conclusion/}{conclusion}
	\import{misc/annexe/}{annexe}
	% And the appendix goes here
	%\appendix
	%\import{chapters/appendices/}{app0A.tex}
	
	%
	% Apparently the guidelines don't say anything about citations or
	% bibliography styles so I guess we can use anything.
	% \backmatter
	% \bibliographystyle{plain}
	% \refstepcounter{chapter}
	% \bibliography{bibliography.bib}
	%Starting of Bibliography %%%%%
	%\begin{thebibliography}{}
	% \bibitem{InvertedPendulum} "Inverted Pendulum", https://en.wikipedia.org/wiki/Inverted\_pendulum
	
	% \bibitem{2} "Application of Kalman Filteing and PID Control for Direct Inverted Pendulum Control ",  http://csuchico-dspace.calstate.edu/bitstream/handle/10211.4/145/4\%2022\%2009\%20Jose\%20Miranda.pdf?sequence=1
	
	% \end{thebibliography}
	
	%
	% Add index
	%\printindex
	%  
%	\part*{Bibliographie}
%	\bibliographystyle{unsrt}
%	\bibliography{bibliography.bib}
\end{document}